% META: {"id":"1SPE-DERLOCAL-CO-012","chapitre":"1SPE-DERIVATION-LOCAL","type_objet":"corrige","exercice_id":"1SPE-DERLOCAL-EX-012","status":"approved","origin":"generated","status_history":["generated","audited","accepted","approved"],"release_acceptance":"RELEASE_OWNER_BATCH_ACCEPTANCE","acceptance_closure_digest":"sha256:9b3ccf9a81c5520fbb7e03b7b2d3e7bf2057a02834b6d908bf3be5f49f3b553f","acceptance_receipt_digest":"sha256:4fad86f0282e0fa4e0419cca1ad6379ae7af5a280c04a4a90880f90a1c044242"}
% BEGIN-VERIFY
% from sympy import symbols, simplify
% h = symbols('h', nonzero=True)
% taux = simplify((1/(1+h) - 1)/h)
% assert simplify(taux + 1/(1+h)) == 0
% END-VERIFY
\begin{corrige}{1SPE-DERLOCAL-EX-012}
\textbf{1.} $f(1+h)-f(1)=\dfrac{1}{1+h}-1=\dfrac{1-(1+h)}{1+h}=\dfrac{-h}{1+h}$.

Donc $\dfrac{f(1+h)-f(1)}{h}=\dfrac{-h}{h(1+h)}=\dfrac{-1}{1+h}$.

L'égalité est bien vérifiée pour $h
e 0$ et $1+h>0$.

\textbf{2.} Lorsque $h$ se rapproche de $0$, $\dfrac{-1}{1+h}$ se rapproche de $\dfrac{-1}{1}=-1$. Donc $f'(1)=-1$.

\textbf{3.} La tangente au point d'abscisse $1$ a pour équation
\[
y=f(1)+f'(1)(x-1)=1+(-1)(x-1)=1-x+1=2-x.
\]

\textbf{4.} La tangente coupe l'axe des abscisses lorsque $y=0$, soit $2-x=0$, c'est-à-dire $x=2$. La tangente coupe l'axe des abscisses au point $(2\,;\,0)$.
\end{corrige}
